\documentclass[11pt]{article}
\usepackage[a4paper,margin=1.05in]{geometry}
\usepackage[T1]{fontenc}
\usepackage[utf8]{inputenc}
\usepackage{lmodern}
\usepackage{amsmath,amssymb,amsthm,mathtools}
\usepackage{enumitem}
\usepackage[hidelinks]{hyperref}
\usepackage{microtype}
\setlength{\emergencystretch}{3em}

\newtheorem{theorem}{Theorem}[section]
\newtheorem{proposition}[theorem]{Proposition}
\newtheorem{lemma}[theorem]{Lemma}
\newtheorem{corollary}[theorem]{Corollary}
\newtheorem{definition}[theorem]{Definition}
\newtheorem{remark}[theorem]{Remark}
\newtheorem{assumption}[theorem]{Assumption}

\title{Predictive Fibres and Critical Oseen Transfer for the Three-Dimensional Navier-Stokes Equations}
\author{Luca Blanchi}
\date{June 2026}

\begin{document}
\maketitle

\begin{abstract}
This article gives a conditional proof strategy for local regularity of the
three-dimensional incompressible Navier-Stokes equations.  The strategy is
presentation-theoretic.  Around a putative singular point, one fixes the
incoming and lateral information visible before the final core of a parabolic
cylinder.  Strong solutions generated from this information are the
predictors.  Suitable weak solutions compatible with the same observations form
the fibre.  A persistent singularity is then represented as a persistent
failure of all admissible predictors to control the fibre across a nested
sequence of critical cylinders.

The main theorem is a conditional reduction.  Given a finite list of
scale-invariant source theorems, it proves that no bad predictive chain can
persist.  The source theorems are stated explicitly in the body of the paper:
they concern the production of the local predictor, the global representation of
residual innovations in one Hilbert space, the packet-level rigidity of nearly
invisible quadratic interactions, the localization of carrier escape, the
synchronization of pressure gauges, and the final predictive
Caffarelli-Kohn-Nirenberg gate.

Several parts of the reduction are proved here in full detail.  These include
the Serrin-critical local Oseen transfer estimate with cutoff and pressure-tail
terms, the relaxed minimax-centre construction, the Hilbertian packing estimate
for genuinely new normal innovations, the two-mode null classification of the
Leray-projected quadratic symbol, the exact pressure absorption of Beltrami
fields, the exact two-and-a-half-dimensional reduction, the relative
Caffarelli-Kohn-Nirenberg gate once predictor smallness is available, the
compactness of the final gate-excess defect, and the fixed-profile cutoff and
fixed-drift subchannels of the Oseen accounting.  The source theorem interface
then records the exact analytic inputs still needed to turn the conditional
reduction into a nonconditional local regularity theorem.
\end{abstract}

\section{Purpose}

The incompressible Navier-Stokes equations in three space dimensions are
\[
\partial_t u-\Delta u+(u\cdot\nabla)u+\nabla p=0,
\qquad
\nabla\cdot u=0.
\]
The global regularity problem remains open. Classical local theory studies suitable weak solutions through scale-invariant quantities, pressure decompositions, local energy inequalities, and epsilon-regularity criteria.

This paper proposes a complementary organization. Around a putative singular
point, one asks whether the core of a parabolic cylinder is causally predictable
from incoming and lateral information by a strong local model. The states
compatible with the same incoming observables form a fibre. A singular chain is
then interpreted as a sequence of scales at which no bounded predictor captures
all states in the fibre.

The presentation-theoretic language is used through analytic statements. The paper separates three levels:
\begin{enumerate}[label=(\roman*)]
\item statements proved here;
\item conditional propositions proved from explicitly stated analytic source theorems;
\item analytic source theorems used as hypotheses in the conditional theorem.
\end{enumerate}
The level of each statement is marked explicitly by its environment: lemma, proposition, theorem, or assumption.

The presentation layer fixes what is visible at a scale, what is invisible
inside the fibre, what it means to update a predictor, and which errors are
allowed to be counted at the next scale.  Standard local regularity estimates
enter as analytic engines inside that organization.  The strategy turns
persistent unpredictability into a finite alternative: either the fibre produces
orthogonal innovations, which are square-summable; or it collapses onto a rigid
branch; or it exits the current presentation through a registered critical
defect.  The conditional theorem below is this alternative written as a scale
iteration.

\section{Main Conditional Statement}

The following theorem is the final form of the reduction proved in this paper.
The detailed definitions are given in the subsequent sections, and the proof is
repeated at the end after all ingredients have been stated.

\begin{theorem}[Conditional predictive reduction]
Let \((u,p)\) be a suitable weak solution in a neighbourhood of \(z_0\).  Assume
 that the local predictive presentations centered at \(z_0\) satisfy the source
theorems stated below: predictor production and differentiability,
global residual representation, packet rigidity, carrier localization, pressure
and gauge synchronization, and final predictive epsilon regularity.  Assume also
that the explicit scale-selection errors are summable.  Then there is no
infinite nested sequence of cylinders at \(z_0\) on which the predictive defect
stays bounded below.  Consequently, if failure of the predictive gate is the only
remaining obstruction to the Caffarelli-Kohn-Nirenberg epsilon condition, then
\(u\) is regular at \(z_0\).
\end{theorem}

The theorem is conditional exactly in the sense expressed by its hypotheses.  All
use of an unproved analytic input occurs through one of those source theorems.
The intervening logic--Oseen transfer, minimax selection, Hilbert packing,
symbolic null classification, rigid branch reduction, compactness of the gate,
and the final contradiction--is proved inside the paper.

\subsection*{Reading guide and dependency map}

The proof has four layers.
\begin{enumerate}[label=(\roman*)]
\item The local analytic layer proves the perturbative estimates used inside one
scale block: the Oseen transfer lemma, fixed-profile cutoff domination,
fixed-drift stability, pressure-tail accounting, relative
Caffarelli-Kohn-Nirenberg gating, and compactness of suitable fibres.
\item The Hilbert layer turns genuinely new linear normal defects into a
square-summable sequence, provided all scale-dependent defects admit a global
residual representation with a Bessel bound.
\item The geometric layer identifies two exact null mechanisms for the projected
quadratic symbol--Beltrami and two-and-a-half-dimensional branches--and records
the general coherent-branch alternative needed beyond pure modes.
\item The source-theorem layer is the conditional interface.  It specifies the
PDE inputs that must be proved to construct predictors, synchronize pressures,
transport defects, handle moving localization, produce branch predictors, and
close the final bad-chain alternative.
\end{enumerate}

Only the fourth layer is assumed.  The first three layers are proved as
standalone statements and then consumed by the modular regularity theorem.

\section{Local Setting}

\subsection{Parabolic scaling}

For \(z_0=(x_0,t_0)\), write
\[
Q_R(z_0)=B_R(x_0)\times(t_0-R^2,t_0).
\]
The Navier-Stokes scaling is
\[
u^{(R)}(x,t)=R u(x_0+Rx,t_0+R^2t),
\qquad
p^{(R)}(x,t)=R^2p(x_0+Rx,t_0+R^2t).
\]
A quantity is critical if its normalized value is stable under this scaling.

\subsection{Suitable weak solutions}

A pair \((u,p)\) is a suitable weak solution on \(Q\) if
\[
u\in L^\infty_tL^2_x(Q)\cap L^2_t\dot H^1_x(Q),
\qquad
p\in L^{3/2}_{t,x}(Q),
\]
the equations hold distributionally, and the local energy inequality holds. For every nonnegative \(\phi\in C_c^\infty(Q)\),
\[
\begin{aligned}
&\int |u(x,t)|^2\phi(x,t)\,dx
+2\int_{-\infty}^t\int |\nabla u|^2\phi\,dx\,ds\\
&\le
\int_{-\infty}^t\int
|u|^2(\partial_s\phi+\Delta\phi)
+(|u|^2+2p)u\cdot\nabla\phi\,dx\,ds.
\end{aligned}
\]

\subsection{Pressure cleaning}

Let \(B\subset\mathbb R^3\) be a smooth ball. For \(1<s<\infty\), the local Helmholtz projection \(\mathbb P_B\) is
\[
\mathbb P_B f=f-\nabla \pi_B,
\]
where \(\pi_B\) solves the Neumann problem
\[
\Delta \pi_B=\nabla\cdot f,
\qquad
\partial_n\pi_B=f\cdot n,
\]
with zero average normalization. Standard elliptic estimates give
\[
\|\mathbb P_B f\|_{L^s(B)}+\|\nabla\pi_B\|_{L^s(B)}
\le C_{s,B}\|f\|_{L^s(B)}.
\]

For a local argument, the pressure is decomposed schematically as
\[
p=p_{\mathrm{loc}}+p_{\mathrm{harm}}+p_{\mathrm{tail}}.
\]
The local part is absorbed by \(\mathbb P_B\). The harmonic and nonlocal pieces contribute a pressure-tail functional
\[
\operatorname{PressTail}(Q_R,Q_r).
\]
In an Oseen estimate with residual \(w\), pressure representative \(\pi\), and
cutoff \(\zeta\), this notation means a bound for the remaining pressure pairing
\[
\left|
\int_{Q_R}\pi_{\mathrm{tail}}\,w\cdot\nabla(\zeta^2)\,dx\,dt
\right|,
\]
or for the corresponding supremum over the admissible energy tests in the same
pressure gauge.  Thus the local projection removes only the local gradient part;
the nonlocal pressure contribution is paid explicitly.

\section{Predictive Presentations}

\subsection{Incoming data}

Fix \(0<\theta<1\) and \(0<\lambda<1\), with \(\theta^2<\lambda\). For \(Q_R(z_0)\), define the incoming region
\[
\mathcal I^-(Q_R)
=
B_R(x_0)\times(t_0-R^2,t_0-\lambda R^2)
\cup
(B_R(x_0)\setminus B_{\theta R}(x_0))\times(t_0-R^2,t_0).
\]
It contains a past slab and a lateral collar, but excludes the final core.

\subsection{Predictors and fibres}

A causal strong predictor on \(Q_R\) is a strong solution \(g\) constructed from data measurable on \(\mathcal I^-(Q_R)\), after the chosen pressure-cleaning procedure. The critical predictor budget is
\[
\|g\|_{L^q_tL^p_x(Q_R)}\le L,
\qquad
\frac2q+\frac3p=1,
\qquad
p>3.
\]
Additional budget components may record pressure-tail size, observable cost, localization cost, and conditioning of the finite observational Gramian.

Let \(\mathcal O\) be a finite family of incoming observables. The invisible fibre associated with \((u,p)\) is the class
\[
\mathsf S^-_{L,\mathcal O}(Q_R)
\]
of suitable weak states \((v,q)\) in \(Q_R\) with budget at most \(L\) and the same incoming observations as \((u,p)\):
\[
\mathcal O(v,q)=\mathcal O(u,p).
\]
Let \(\mathsf P^-_{L,\mathcal O}(Q_R)\) denote the corresponding class of causal strong predictors.

\begin{definition}[Local predictive defect]
Let \(d_{Q_R}\) be a pressure-cleaned distance on the core \(Q_{\theta R}\). The minimax predictive defect is
\[
\mathfrak D_{L,\mathcal O}(Q_R)
=
\inf_{g\in \mathsf P^-_{L,\mathcal O}(Q_R)}
\sup_{(v,q)\in \mathsf S^-_{L,\mathcal O}(Q_R)}
d_{Q_R}((v,q),g).
\]
When the infimum is taken in the closed convex hull of the predictor class, the distance from the relaxed centre to realizable predictors is recorded as a realization defect \(\mathfrak V(Q_R)\).
\end{definition}

A bad predictive chain at \(z_0\) is a nested sequence \(Q_{R_k}(z_0)\), \(R_k\downarrow0\), for which \(\mathfrak D_{L,\mathcal O_k}(Q_{R_k})\) stays bounded below by a fixed positive number.

\section{Critical Local Oseen Transfer}

The following estimate is the local perturbative engine.  It transports a
residual around a Serrin-critical strong predictor from a cylinder to a smaller
core, while recording the cutoff, drift, forcing, and pressure-tail costs that
must be paid by the source theorem interface.

\begin{lemma}[Critical local Oseen transfer]
Let \(Q_\rho\Subset Q_R\). Let \(g\) be divergence-free and
\[
g\in L^q_tL^p_x(Q_R),
\qquad
\frac2q+\frac3p=1,
\qquad
p>3.
\]
Let \(w\) be divergence-free and pressure-cleaned, satisfying
\[
\partial_t w-\Delta w
+(g\cdot\nabla)w
+(w\cdot\nabla)g
+\nabla \pi
=F
\]
in \(Q_R\), distributionally. Choose \(\zeta\in C_c^\infty(Q_R)\) with
\(\zeta=1\) on \(Q_\rho\). Assume that the nonlocal pressure contribution is
bounded by \(\operatorname{PressTail}_{\zeta}\), in the sense of the pressure
tail pairing defined above, and define
\[
\begin{aligned}
\operatorname{Coll}_{\zeta,g}(w)^2
=&
\sup_{\mathrm{initial\ slab}}\int |w|^2\zeta^2\,dx\\
&+
\int_{Q_R}|w|^2
\left(
|\partial_t\zeta|+|\Delta\zeta|+|\nabla\zeta|^2
\right)\,dx\,dt\\
&+
\int_{Q_R}|g|\,|w|^2|\nabla\zeta|\,dx\,dt.
\end{aligned}
\]
Then
\[
\|w\|_{L^\infty_tL^2_x(Q_\rho)}
+\|\nabla w\|_{L^2_{t,x}(Q_\rho)}
+\|w\|_{L^{10/3}_{t,x}(Q_\rho)}
\]
\[
\le
C
\left(
\operatorname{Coll}_{\zeta,g}(w)
+\|F\|_{L^2_tH^{-1}_x(Q_R)}
+\operatorname{PressTail}_{\zeta}
\right),
\]
where \(C\) depends on \(Q_\rho\Subset Q_R\), on \(p\), and on
\(\|g\|_{L^q_tL^p_x(Q_R)}\).
\end{lemma}

\begin{proof}
Test the equation against \(w\zeta^2\).  The local pressure projection removes
the local pressure gradient.  The remaining pressure pairing is bounded by
\(\operatorname{PressTail}_{\zeta}\).  We obtain
\[
\frac12\frac{d}{dt}\int |w|^2\zeta^2
+\int |\nabla w|^2\zeta^2
\le
I_{\mathrm{cut}}+I_g+I_F+I_{\mathrm{press}}.
\]
The cutoff term is bounded by the collar part of
\(\operatorname{Coll}_{\zeta,g}(w)^2\).  The forcing term is estimated by
duality:
\[
|\langle F,w\zeta^2\rangle|
\le
\|F\|_{H^{-1}}\|w\zeta^2\|_{H^1}
\le
\frac14\|\zeta\nabla w\|_2^2
+C\|F\|_{H^{-1}}^2
+C\|w\nabla\zeta\|_2^2.
\]
Since \(\nabla\cdot g=0\),
\[
\int (g\cdot\nabla w)\cdot w\zeta^2
=
-\frac12\int |w|^2g\cdot\nabla(\zeta^2),
\]
which is included in the collar cost.

It remains to estimate the stretching term. Using \(\nabla\cdot w=0\),
\[
\int w_i(\partial_i g_j)w_j\zeta^2
=
-\int g_jw_i(\partial_i w_j)\zeta^2
-\int g_jw_iw_j\partial_i(\zeta^2).
\]
The second term is again a collar term. For the first,
\[
\left|
\int g_jw_i(\partial_i w_j)\zeta^2
\right|
\le
\|g\|_{L^p}
\|w\zeta\|_{L^{2p/(p-2)}}
\|\zeta\nabla w\|_{L^2}.
\]
By Gagliardo-Nirenberg,
\[
\|w\zeta\|_{L^{2p/(p-2)}}
\le
C
\|w\zeta\|_{L^2}^{1-3/p}
\|\nabla(w\zeta)\|_{L^2}^{3/p}.
\]
Writing
\[
E(t)=\int |w|^2\zeta^2,
\qquad
D(t)=\int |\nabla w|^2\zeta^2,
\]
and absorbing the cutoff contribution into the collar cost, Young's inequality gives
\[
|I_g^{\mathrm{int}}|
\le
\frac14D(t)+C\|g(t)\|_{L^p}^{q}E(t),
\qquad
q=\frac{2p}{p-3}.
\]
The exponent is exactly the Serrin exponent because \(2/q+3/p=1\).

Thus
\[
\frac{d}{dt}E(t)+\frac12D(t)
\le
C\|g(t)\|_{L^p}^{q}E(t)+B(t),
\]
where \(B(t)\) contains only forcing, pressure-tail, and collar terms. Since \(\|g(t)\|_{L^p}^{q}\in L^1_t\), Gronwall yields
\[
\sup_tE(t)+\int D(t)\,dt
\le
C
\left(
\operatorname{Coll}_{\zeta,g}(w)^2
+\|F\|_{L^2_tH^{-1}_x(Q_R)}^2
+\operatorname{PressTail}_{\zeta}^2
\right).
\]
Finally,
\[
\|w\|_{L^{10/3}(Q_\rho)}
\lesssim
\|w\|_{L^\infty_tL^2_x(Q_\rho)}^{2/5}
\|\nabla w\|_{L^2(Q_\rho)}^{3/5}.
\]
This proves the claim.
\end{proof}

\subsection{Fixed-profile collars and the pressure interface}

The Oseen estimate becomes a scale iteration only after the cutoff family has
been fixed.  The useful closed subchannel is the integrated derivative-weighted
cutoff collar for a single anchored dyadic family.

\begin{proposition}[Fixed-profile dyadic cutoff domination]
Fix \(0<a<b<1\), \(0<\lambda<1\), a centre \(z_*=(x_*,t_*)\), and
\(r_k=2^{-k}R\).  Let
\[
\zeta_k(x,t)
=
\phi\!\left(\frac{x-x_*}{r_k}\right)
\psi\!\left(\frac{t_*-t}{r_k^2}\right),
\]
where \(\phi\in C_c^\infty(B_b(0))\), \(\phi=1\) on \(B_a(0)\), and \(\psi\)
is a fixed smooth past-time cutoff.  Put
\[
W_k
=
|\partial_t\zeta_k|+|\Delta\zeta_k|+|\nabla\zeta_k|^2 .
\]
Then, inside the fixed outer cylinder,
\[
\sum_{k\ge K} W_k(x,t)
\le
C\left(|x-x_*|^{-2}+R^{-2}\right)
\]
for a constant depending only on the profiles and on \(a,b,\lambda\).  Hence, for
every local energy field \(w\),
\[
\sum_{k\ge K}\int |w|^2 W_k
\le
C
\left(
\int |\nabla w|^2
+R^{-2}\int |w|^2
\right)
\]
on the outer cylinder, after the standard local Hardy inequality.
\end{proposition}

\begin{proof}
The derivative bounds give
\[
W_k\le
C r_k^{-2}
\left(
{\bf 1}_{A_k^{\mathrm{lat}}}
+{\bf 1}_{A_k^{\mathrm{time}}}
\right),
\]
where \(A_k^{\mathrm{lat}}\) is the spatial annulus
\(a r_k\le |x-x_*|\le b r_k\) during the active time interval, and
\(A_k^{\mathrm{time}}\) is the past time floor on which the time cutoff varies.
On \(A_k^{\mathrm{lat}}\), \(r_k\) is comparable to \(|x-x_*|\), and a fixed point
belongs to only boundedly many dyadic annuli.  Thus the sum of the lateral
weights is bounded by \(C|x-x_*|^{-2}\).

On \(A_k^{\mathrm{time}}\), the condition
\(t_*-t\simeq r_k^2\) allows only boundedly many dyadic indices for each fixed
time.  For those indices, the spatial support lies in \(|x-x_*|\le b r_k\);
the same dyadic comparison gives the pointwise bound away from \(x=x_*\), which
is sufficient for the integrated estimate.  A possible outer cutoff contributes
only \(C R^{-2}\).  Integrating against \(|w|^2\) and applying the local Hardy
inequality gives the displayed bound.
\end{proof}

\begin{proposition}[Fixed-drift collar stability]
Let \(g\) be fixed on a block of nested cylinders and suppose
\[
\|g\|_{L^q_tL^p_x(Q_R)}\le L,
\qquad
\frac2q+\frac3p=1,
\qquad
p>3.
\]
For every perturbation \(w\) satisfying the local Oseen equation on the block,
the drift-dependent part of the collar is controlled by the same annular
perturbation norm used in the Oseen transfer:
\[
\int_{Q_R}|g|\,|w|^2|\nabla\zeta_R|
\le
C(L,\theta,p)
\|w\|_{L^\infty_tL^2_x(Q_R\setminus Q_{\theta R})}^{2(1-3/p)}
\|\nabla w\|_{L^2(Q_R\setminus Q_{\theta R})}^{6/p}.
\]
Thus a fixed drift does not introduce a new scale owner: it is absorbed by the
same annular energy and dissipation already paid in the Oseen estimate.
\end{proposition}

\begin{proof}
Use Hölder exactly as in the preceding proposition and interpolate
\[
\|w\|_{L^{2p/(p-2)}_x}
\le
C
\|w\|_{L^2_x}^{1-3/p}
\|\nabla w\|_{L^2_x}^{3/p}
\]
on each time slice.  Integrating in time with the Serrin conjugate exponent gives
the claimed bound.  Since \(g\) and the cutoff profile are fixed before the block
is selected, the term is a predetermined Oseen collar contribution.
\end{proof}

\begin{proposition}[Pressure-pair source interface]
For a fixed dyadic Oseen chart, the pressure pairing
\[
\operatorname{PressPair}_k
=
\left|\int \pi_k^0\,w_k\cdot\nabla(\zeta_k^2)\right|
\]
reduces, after parabolic normalization and Young's inequality, to critical
pressure and velocity collar densities:
\[
\operatorname{PressPair}_k
\le
\varepsilon V_k+C_\varepsilon P_k,
\]
where
\[
V_k=r_k^{-2}\int_{A_k}|w_k|^3,
\qquad
P_k=r_k^{-2}\int_{A_k}|\pi_k^0|^{3/2}.
\]
The fixed-profile Hardy estimate above does not by itself imply summability of
\(\sum_k P_k\) or \(\sum_k V_k\).  Those two normalized collar sums are part of
the pressure/localization source theorem interface.
\end{proposition}

\begin{proof}
The first displayed inequality is the usual local pressure estimate on the
normalized chart, written back in scale-invariant variables.  It separates the
pressure pairing into a velocity collar and a pressure collar.

The obstruction is purely scaling.  Let \(A_k\) be disjoint dyadic annuli and
choose a nonnegative density \(m\) with
\[
\int_{A_k}m=r_k^2/k.
\]
Then \(\sum_k\int_{A_k}m<\infty\), while
\[
\sum_k r_k^{-2}\int_{A_k}m
=
\sum_k\frac1k
=\infty .
\]
Taking \(m=|\pi|^{3/2}\), or \(m=|w|^3\), shows that finite local critical
integrability does not produce normalized dyadic collar summability.  Therefore
the pressure-pair term cannot be closed by the fixed cutoff Hardy argument
alone; it needs an independent pressure/localization theorem, a gate, or a
registered source alternative.
\end{proof}

\section{Minimax Centres}

The minimax construction used by the predictive framework is a Hilbert-space statement. The PDE content lies in verifying that the chosen predictor classes satisfy the hypotheses.

\begin{lemma}[Relaxed minimax centres]
Let \(\mathcal H\) be a Hilbert space. Let \(P\subset\mathcal H\) be nonempty, bounded, closed, and convex, and let \(S\subset\mathcal H\) be bounded. Define
\[
R(g)=\sup_{v\in S}\|v-g\|_{\mathcal H},
\qquad
g\in P.
\]
Then \(R\) attains its minimum on \(P\).
\end{lemma}

\begin{proof}
The map \(g\mapsto\|v-g\|_{\mathcal H}\) is convex and weakly lower semicontinuous for each \(v\). Hence \(R\), as a supremum of weakly lower semicontinuous functions, is weakly lower semicontinuous. Since Hilbert spaces are reflexive, a bounded closed convex subset is weakly compact. Therefore \(R\) attains its minimum on \(P\).
\end{proof}

\begin{remark}
In the Navier-Stokes application, \(P\) is the weakly closed convex hull of admissible strong predictors at a fixed finite observation budget. The gap between a relaxed minimax centre and an actual strong predictor is exactly the realization defect \(\mathfrak V\). This defect must be paid for in the scale-transfer inequality.
\end{remark}

\section{Innovation Packing}

\subsection{Abstract packing}

Let \(\mathfrak H\) be a Hilbert space of normalized adjoint observables. Suppose that a sequence of normal defects \(\mathcal N_k\in\mathfrak H^*\) has normalized innovative witnesses \(e_k\in\mathfrak H\). The filtration condition is
\[
e_k\perp \mathcal P_{k-1},
\qquad
\mathcal P_k=\overline{\mathcal P_{k-1}+\operatorname{span}\{e_k\}},
\]
so \((e_k)\) is orthonormal on innovative scales.

\begin{lemma}[Hilbertian innovation packing]
Assume that there exist \(\mathcal R\in\mathfrak H^*\), errors \(\rho_k\), and \(c>0\) such that
\[
\langle \mathcal N_k,e_k\rangle
=
\langle \mathcal R,e_k\rangle+\rho_k,
\qquad
\sum_k|\rho_k|^2<\infty,
\]
and
\[
|\langle \mathcal N_k,e_k\rangle|\ge c\|\mathcal N_k\|_{\mathfrak H^*}.
\]
Then
\[
\sum_k\|\mathcal N_k\|_{\mathfrak H^*}^2
\le
C_c\left(\|\mathcal R\|_{\mathfrak H^*}^2+\sum_k|\rho_k|^2\right).
\]
\end{lemma}

\begin{proof}
By Bessel's inequality,
\[
\sum_k|\langle\mathcal R,e_k\rangle|^2
\le
\|\mathcal R\|_{\mathfrak H^*}^2.
\]
Using the coefficient representation,
\[
\sum_k|\langle\mathcal N_k,e_k\rangle|^2
\le
2\sum_k|\langle\mathcal R,e_k\rangle|^2
+2\sum_k|\rho_k|^2.
\]
The witness lower bound then gives the stated estimate.
\end{proof}

\begin{proposition}[Bessel representation with moving charts]
Let \(H_{\mathrm{def}}\) be a Hilbert space of registered defect vectors.  For
each scale \(k\), let \(d_k\in H_{\mathrm{def}}\), let
\[
i_k=d_k-P_{\overline{\operatorname{span}\{d_j:j<k\}}}d_k,
\]
and, when \(i_k\ne0\), set \(e_k=i_k/\|i_k\|\).  Suppose that there are a Hilbert
space \(\mathcal K\), a residual vector \(\mathcal R\in\mathcal K\), bounded
operators \(A_k:\mathcal K\to H_{\mathrm{def}}\), and errors \(\rho_k\) such that
\[
d_k=A_k\mathcal R+\rho_k,
\qquad
\sum_k\|\rho_k\|_{H_{\mathrm{def}}}^2<\infty,
\]
and such that the pulled-back innovation tests satisfy the Bessel bound
\[
\sum_k|\langle f,A_k^*e_k\rangle_{\mathcal K}|^2
\le
B\|f\|_{\mathcal K}^2
\qquad
\text{for every } f\in\mathcal K .
\]
Then
\[
\sum_k\|i_k\|_{H_{\mathrm{def}}}^2
\le
2B\|\mathcal R\|_{\mathcal K}^2
+2\sum_k\|\rho_k\|_{H_{\mathrm{def}}}^2 .
\]
\end{proposition}

\begin{proof}
For nonzero \(i_k\), the vector \(e_k\) is orthogonal to all previous \(d_j\), so
\[
\|i_k\|
=
\langle d_k,e_k\rangle
=
\langle\mathcal R,A_k^*e_k\rangle+\langle\rho_k,e_k\rangle .
\]
After squaring, summing, and using \(\|e_k\|=1\), the Bessel bound gives the
displayed estimate.  Zero-innovation indices contribute nothing.
\end{proof}

\subsection{PDE meaning}

The abstract packing lemma applies when the transported defect vectors already
live in one fixed Hilbert space.  The moving-chart proposition is the form needed
for Navier-Stokes: scale changes, pressure gauges, carrier frames, and cutoffs
may vary with \(k\), and the maps \(A_k\) must satisfy a genuine Bessel or
finite-overlap estimate.  The PDE source theorem supplies exactly this
representation and the square-summability of the error vectors \(\rho_k\).

\section{The Bilinear Symbol and Null Fibres}

The next lemma is a finite-dimensional calculation for the Leray-projected quadratic symbol. It explains why large invisible fibres with very small quadratic leakage should be close to special geometric branches, but by itself it treats only pure modes.

\begin{lemma}[Two-mode null classification]
Let
\[
u=a e^{i\xi\cdot x},
\qquad
v=b e^{i\eta\cdot x},
\]
with
\[
a\cdot\xi=0,
\qquad
b\cdot\eta=0.
\]
Set \(k=\xi+\eta\), \(k\ne0\), and consider the symmetrized projected symbol
\[
\mathfrak b((\xi,a),(\eta,b))
=
\mathbb P_k\left[(a\cdot\eta)b+(b\cdot\xi)a\right].
\]
Write
\[
a=\alpha k+A,
\qquad
b=\beta k+B,
\qquad
A\perp k,
\qquad
B\perp k.
\]
Then \(\mathfrak b=0\) if and only if
\[
\alpha B+\beta A=0.
\]
The nonzero solutions fall into the following branches:
\begin{enumerate}[label=(\roman*)]
\item the shear-null branch \(\alpha=\beta=0\);
\item the balanced equal-radius branch \(\alpha\beta\ne0\), which forces \(|\xi|=|\eta|\).
\end{enumerate}
\end{lemma}

\begin{proof}
Since \(a\cdot\xi=0\) and \(b\cdot\eta=0\),
\[
a\cdot\eta=a\cdot(\xi+\eta)=a\cdot k,
\qquad
b\cdot\xi=b\cdot(\xi+\eta)=b\cdot k.
\]
Therefore
\[
(a\cdot\eta)b+(b\cdot\xi)a
=
(a\cdot k)b+(b\cdot k)a
=
|k|^2(\alpha b+\beta a).
\]
Substituting the decompositions gives
\[
\alpha b+\beta a
=
2\alpha\beta k+\alpha B+\beta A.
\]
The projection \(\mathbb P_k\) kills the component parallel to \(k\), so the projected symbol vanishes exactly when \(\alpha B+\beta A=0\).

If \(\alpha=\beta=0\), both polarizations are perpendicular to \(k\), giving the shear-null branch. If exactly one of \(\alpha,\beta\) is zero, then the relation forces one polarization to vanish, so there is no nonzero two-mode interaction. If \(\alpha\beta\ne0\), then \(B=-(\beta/\alpha)A\). The incompressibility condition for \(a\) gives
\[
A\cdot\xi=-\alpha k\cdot\xi.
\]
Since \(A\perp k\), we have \(A\cdot\eta=-A\cdot\xi=\alpha k\cdot\xi\). The incompressibility condition for \(b\) gives
\[
0=b\cdot\eta=\beta k\cdot\eta+B\cdot\eta
=
\beta k\cdot\eta-\frac{\beta}{\alpha}A\cdot\eta
=
\beta(k\cdot\eta-k\cdot\xi).
\]
Thus \(k\cdot\eta=k\cdot\xi\), which is
\[
|\eta|^2+\xi\cdot\eta=|\xi|^2+\xi\cdot\eta.
\]
Hence \(|\xi|=|\eta|\).
\end{proof}

\begin{remark}
The next analytic step is a quantitative packet-level null-graph rigidity theorem: it should promote this exact two-mode calculation to localized parabolic wave packets and identify when a large fibre with small leakage is close to a Beltrami, shear, or carrier-escape branch.
\end{remark}

\section{Rigid Branches}

\begin{proposition}[Beltrami absorption]
Suppose \(U\) is divergence-free and
\[
\nabla\times U=\lambda U.
\]
Then
\[
(U\cdot\nabla)U
=
\nabla\left(\frac{|U|^2}{2}\right).
\]
Thus the quadratic term is pure pressure and is eliminated by Leray projection or by local pressure-cleaning.
\end{proposition}

\begin{proof}
Use the identity
\[
(U\cdot\nabla)U
=
\nabla\left(\frac{|U|^2}{2}\right)
-U\times(\nabla\times U).
\]
If \(\nabla\times U=\lambda U\), then \(U\times(\lambda U)=0\).
\end{proof}

\begin{proposition}[Two-and-a-half-dimensional reduction]
Let
\[
U(x,t)=V(y,t)+W(y,t)n,
\]
where \(y\) belongs to a fixed plane \(\Pi\), \(V(y,t)\in\Pi\), \(n\perp\Pi\), and \(\nabla_\Pi\cdot V=0\). Then the three-dimensional Navier-Stokes system reduces to
\[
\partial_t V+V\cdot\nabla_\Pi V-\Delta_\Pi V+\nabla_\Pi P=0,
\]
\[
\partial_t W+V\cdot\nabla_\Pi W-\Delta_\Pi W=0.
\]
Thus the planar part is a two-dimensional Navier-Stokes solution and the transverse component is a passive transported-diffused scalar.
\end{proposition}

\begin{proof}
All fields are constant in the \(n\)-direction. The in-plane component of the equation is exactly the two-dimensional incompressible Navier-Stokes equation for \(V\). The \(n\)-component contains no pressure gradient and gives the passive scalar equation for \(W\).
\end{proof}

\section{Relative Caffarelli-Kohn-Nirenberg Gates}

The Caffarelli-Kohn-Nirenberg criterion is used here in a relative form.  The
predictor is allowed to carry the large, but controlled, Serrin component; the
unknown part is the perturbation around that predictor.

For a cylinder \(Q_R(z_0)\), write
\[
C_R(u,p;z_0)
=
R^{-2}\int_{Q_R(z_0)}|u|^3\,dx\,dt
+
R^{-2}\int_{Q_R(z_0)}
|p-(p)_{B_R}(t)|^{3/2}\,dx\,dt.
\]

\begin{proposition}[Relative Caffarelli-Kohn-Nirenberg criterion]
Let \(g\) be a strong divergence-free solution on \(Q_1\) with pressure \(\pi_g\),
and suppose
\[
\|g\|_{L^q_tL^p_x(Q_1)}\le L,
\qquad
\frac2q+\frac3p=1,
\qquad
p>3.
\]
Let \((u,p)\) be suitable on \(Q_1\), set \(w=u-g\), and write
\[
p-\pi_g=\pi_w+\pi_{\mathrm{tail}}
\]
with the local pressure part normalized by zero spatial mean on balls.  For every
\(\varepsilon_{\mathrm{CKN}}>0\) there exist
\(\theta=\theta(L,p,\varepsilon_{\mathrm{CKN}})\in(0,1)\) and
\(\delta=\delta(L,p,\varepsilon_{\mathrm{CKN}})>0\) such that, if
\[
\|w\|_{L^{10/3}(Q_1)}
+
\|\pi_w\|_{L^{3/2}(Q_1)}
+
\|\pi_{\mathrm{tail}}\|_{L^{3/2}(Q_1)}
\le \delta,
\]
then
\[
C_\theta(u,p;0)\le \varepsilon_{\mathrm{CKN}}.
\]
In particular, choosing \(\varepsilon_{\mathrm{CKN}}\) below the classical
Caffarelli-Kohn-Nirenberg threshold gives regularity in a smaller cylinder.
\end{proposition}

\begin{proof}
By Hölder and the Serrin bound, \(g\in L^3(Q_\theta)\) with
\[
\theta^{-2}\int_{Q_\theta}|g|^3
\le
C(L,p)\theta^\alpha
\]
for some \(\alpha>0\), because \(p>3\) gives subcritical local integrability on
shrinking cylinders after the Serrin normalization.  The associated pressure
\(\pi_g\) obeys the same shrinking estimate after subtracting its spatial mean,
by local Calderon-Zygmund estimates and the harmonic-pressure decomposition.
Choose \(\theta\) so these two predictor contributions are at most
\(\varepsilon_{\mathrm{CKN}}/4\).

For the perturbation, interpolation gives
\[
\theta^{-2}\int_{Q_\theta}|w|^3
\le
C\|w\|_{L^{10/3}(Q_1)}^3\theta^{\beta-2}
\]
with a fixed positive exponent \(\beta\) after the parabolic volume factor is
accounted for.  The pressure perturbation contributes
\[
\theta^{-2}\int_{Q_\theta}|\pi_w|^{3/2}
\le
\theta^{-2}\|\pi_w\|_{L^{3/2}(Q_1)}^{3/2},
\]
and the same estimate applies to \(\pi_{\mathrm{tail}}\).  With \(\theta\) fixed,
choose \(\delta\) so the three perturbative contributions are together at most
\(\varepsilon_{\mathrm{CKN}}/2\).  Summing predictor and perturbation pieces gives
the displayed bound for \(C_\theta(u,p;0)\).
\end{proof}

\subsection{Final gate-excess defect}

The presentation uses a final gate to decide whether a scale is already within
the relative Caffarelli-Kohn-Nirenberg regime.  The following compactness fact is
the precise topological statement needed in the iteration.

\begin{definition}[Gate-excess defect]
Fix a bounded presentation chart, a predictor class \(P\), a suitable weak fibre
\(S\), a core cylinder \(Q_\theta\), and a gate threshold
\(\varepsilon_{\mathrm{gate}}>0\).  For \(g\in P\) and \((v,q)\in S\), define
\[
\Gamma((v,q),g)
=
\left[
C_\theta(v,q\mid g)-\varepsilon_{\mathrm{gate}}
\right]_+,
\]
where \(C_\theta(v,q\mid g)\) denotes the relative Caffarelli-Kohn-Nirenberg
quantity of \(v-g\) with the corresponding pressure difference and pressure-tail
allowance.  The gate-excess defect of the chart is
\[
\mathfrak G
=
\inf_{g\in P}\sup_{(v,q)\in S}\Gamma((v,q),g).
\]
\end{definition}

\begin{proposition}[Gate-excess dichotomy]
Assume \(P\) is compact in the strong topology used for predictors, \(S\) is
compact in the suitable-weak topology, and the relative Caffarelli-Kohn-Nirenberg
functional is lower semicontinuous on \(S\times P\).  Then:
\begin{enumerate}[label=(\roman*)]
\item \(\mathfrak G=0\) if and only if every state in the fibre is captured by
some predictor at the final gate;
\item if the final gate is closed and no predictor captures the whole fibre, then
\(\mathfrak G>0\);
\item for each \(\delta>0\), the condition \(\mathfrak G\ge\delta\) is stable
under compact convergence of charts.
\end{enumerate}
\end{proposition}

\begin{proof}
The map
\[
g\mapsto \sup_{(v,q)\in S}\Gamma((v,q),g)
\]
is lower semicontinuous as a supremum of lower semicontinuous functions.  Since
\(P\) is compact, the infimum is attained.  If the attained value is zero, the
chosen predictor has zero gate excess against every fibre state, which is exactly
capture by the final gate.  Conversely, a capturing predictor gives zero excess.
If no such predictor exists, the attained value cannot be zero; hence it is
strictly positive.  Finally, lower semicontinuity under compact convergence gives
closedness of the superlevel sets \(\{\mathfrak G\ge\delta\}\).
\end{proof}

\section{Compactness and Zero-Cost Branches}

The conditional proof repeatedly passes to limits of bounded scale-normalized
objects.  The following propositions isolate what is used and what is not.

\begin{proposition}[Suitable limit passage]
Let \((u_n,p_n)\) be suitable weak solutions on \(Q_1\) with uniform bounds
\[
\sup_n
\left(
\|u_n\|_{L^\infty_tL^2_x(Q_1)}
+\|\nabla u_n\|_{L^2(Q_1)}
+\|p_n\|_{L^{3/2}(Q_1)}
\right)<\infty.
\]
Then a subsequence converges to a suitable weak solution \((u,p)\) on compact
subcylinders of \(Q_1\), with strong convergence of \(u_n\) in \(L^2_{\mathrm{loc}}\)
and lower semicontinuity of the local energy inequality.
\end{proposition}

\begin{proof}
The equation gives a uniform bound for \(\partial_tu_n\) in a negative Sobolev
space on compact subcylinders.  Aubin-Lions compactness gives strong
\(L^2_{\mathrm{loc}}\) convergence of \(u_n\), while the pressure bound gives weak
\(L^{3/2}_{\mathrm{loc}}\) convergence of \(p_n\).  The nonlinear term passes to
the limit because \(u_n\to u\) strongly in \(L^2_{\mathrm{loc}}\) and
\(\nabla u_n\rightharpoonup \nabla u\) weakly in \(L^2_{\mathrm{loc}}\).  The
local energy inequality passes by weak lower semicontinuity of the dissipation
and strong convergence of the lower-order terms.
\end{proof}

\begin{proposition}[Normal-defect limit passage]
Suppose a sequence of normalized presentations has uniformly bounded predictors,
uniformly bounded suitable fibres, convergent pressure gauges, and normal defect
functionals \(\mathcal N_n\) converging weakly to \(\mathcal N\) in a common dual
space.  If the corresponding witness spaces converge by compact embeddings and
the witness lower bounds are uniform, then every nonzero limiting normal defect
is detected by a limiting witness.  In particular, a vanishing witness pairing
for all limiting witnesses forces \(\mathcal N=0\).
\end{proposition}

\begin{proof}
Let \(e\) be a unit witness for the limit space.  By the compact convergence of
witness spaces, choose witnesses \(e_n\) in the \(n\)-th space with \(e_n\to e\)
strongly.  Weak convergence of \(\mathcal N_n\) and strong convergence of
witnesses give
\[
\langle \mathcal N_n,e_n\rangle\to \langle \mathcal N,e\rangle.
\]
If every limiting pairing is zero, then \(\mathcal N\) annihilates the closure of
the witness span.  The uniform witness lower bound says this span is norming for
normal defects, hence \(\mathcal N=0\).
\end{proof}

\begin{proposition}[Euler-coherent pressure branch]
Let \(U\) be a smooth divergence-free field on a cylinder and assume its
quadratic nonlinearity is a pure pressure gradient:
\[
\mathbb P\bigl((U\cdot\nabla)U\bigr)=0.
\]
Then the field \(U\) is invisible to the Leray-projected quadratic normal defect.
Consequently, a zero-cost branch in the predictive presentation must include not
only Beltrami and two-and-a-half-dimensional examples, but every coherent branch
for which the projected quadratic term vanishes after the selected pressure
gauge.
\end{proposition}

\begin{proof}
The Leray projection annihilates gradients.  If
\((U\cdot\nabla)U=\nabla P_U\) in the chosen local gauge, then the projected
quadratic contribution is zero.  Therefore no normal witness built only from the
projected quadratic leakage can distinguish such a branch.  This proves the
claim and explains why the branch alternative in the conditional theorem is
formulated as a general coherent zero-cost branch rather than only the two
explicit model branches proved above.
\end{proof}

\section{Branch Gates and Nonlinear Witnesses}

The branch alternative becomes useful only when it feeds the final gate rather
than merely explaining a loss of quadratic leakage.  The next statement records
the exact implication needed from a branch construction.

\begin{proposition}[Branch gate criterion]
Let a coherent branch predictor \(G\) be produced on a cylinder \(Q_R\), with
branch pressure \(P_G\), and suppose:
\begin{enumerate}[label=(\roman*)]
\item the predictor has a Serrin bound
\(\|G\|_{L^q_tL^p_x(Q_R)}\le L\), \(2/q+3/p=1\), \(p>3\);
\item the relative perturbation obeys
\[
\|u-G\|_{L^{10/3}(Q_R)}
+
\|p-P_G\|_{L^{3/2}(Q_R)}
+
\operatorname{PressTail}(Q_R,Q_{\theta R})
\le\delta;
\]
\item all pressure gauges used to define \(P_G\) and \(p\) are synchronized on
the smaller cylinder.
\end{enumerate}
Then, for \(\delta\) small depending only on \(L,p,\theta\), the final relative
Caffarelli-Kohn-Nirenberg gate opens on \(Q_{\theta R}\).
\end{proposition}

\begin{proof}
After parabolic normalization, this is exactly the relative
Caffarelli-Kohn-Nirenberg criterion.  The pressure synchronization hypothesis is
needed only to ensure that the pressure difference entering the relative
quantity is the same pressure difference controlled in the branch construction.
With the gauges aligned, the smallness assumption gives the perturbative
quantity required by the criterion.
\end{proof}

\begin{proposition}[Nonlinear Caffarelli-Kohn-Nirenberg witnesses are not
Hilbert innovations]
The raw Caffarelli-Kohn-Nirenberg quantity
\[
C_R(u,p;z_0)
=R^{-2}\int_{Q_R}|u|^3
+R^{-2}\int_{Q_R}|p-(p)_{B_R}(t)|^{3/2}
\]
cannot, by itself, be inserted as a linear Hilbert innovation in the Bessel
packing argument.  Any use of this quantity in the predictive proof must be
routed through one of the following mathematical alternatives: a relative gate,
a linearized source functional with a norming witness, a compactness-rigidity
branch, or a separate monotone/Carleson estimate.
\end{proposition}

\begin{proof}
The Hilbert packing lemma applies to linear functionals paired with orthonormal
witnesses.  The map \(u\mapsto\int |u|^3\) is nonlinear and has a large kernel
under linearization at \(u=0\); the pressure part is also nonlinear through the
quadratic pressure relation.  Therefore raw Caffarelli-Kohn-Nirenberg mass is not
a Hilbert coefficient.  To use it in a square-summability argument one must first
produce a genuine linear defect functional, or else use the quantity as a gate
whose success exits the iteration.  The listed alternatives are exactly those
two possibilities and their compactness variants.
\end{proof}

\begin{proposition}[Gated witnesses exit the source count]
If a Caffarelli-Kohn-Nirenberg witness opens the final relative gate at a scale,
then that scale is terminal for the bad-chain argument and its witness mass is
not counted in the Hilbert source budget.  If the gate does not open, the witness
may be used only after it has been converted into a source functional, a rigidity
alternative, or an escape alternative.
\end{proposition}

\begin{proof}
Opening the final gate gives regularity in a smaller cylinder by the relative
criterion and the classical epsilon theorem.  A bad chain cannot continue through
that cylinder, so the scale is no longer part of the infinite source count.  If
the gate remains closed, no regularity conclusion has been obtained; counting the
raw nonlinear witness as a Hilbert coefficient would be exactly the invalid step
excluded by the preceding proposition.  Thus the two cases are disjoint.
\end{proof}

\section{Source Theorem Interface}

We now state the analytic inputs consumed by the reduction.  Each source theorem
has a specific output: an object, a summability statement, a branch, a gate, or a
registered defect.  The modular proof never uses an unnamed analytic input.

\begin{proposition}[Finite presentation normal form]
Assume that along a tail of a predictive chain the presentation alphabet is
fixed, the admissible transition rule between consecutive scale states is total,
and every accepted positive component is assigned to exactly one source ledger.
Then every transition on that tail has one of the following statuses:
\begin{enumerate}[label=(\roman*)]
\item accepted source term;
\item finite update;
\item relative gate;
\item registered defect;
\item coherent branch;
\item carrier recentering;
\item impossible branch.
\end{enumerate}
Consequently, on a no-gate, no-defect tail with only finitely many updates, the
modular proof may consume only the named source ledgers and the branch or carrier
alternatives stated below.
\end{proposition}

\begin{proof}
This is a finite normal-form argument.  The fixed alphabet prevents new source
names or new gauges from appearing after the tail has been selected.  Totality of
the transition rule classifies each adjacent scale transition.  The owner map
assigns every accepted positive component to one source ledger, so no accepted
cost can be used twice or disappear.  After finitely many finite updates are
discarded, the listed statuses exhaust all transitions on the tail.
\end{proof}

\begin{assumption}[Predictive epsilon regularity source theorem]
For every \(L\) and every admissible pressure-cleaning scheme, there exist
\(\varepsilon(L)>0\) and \(0<\theta<1\) such that the following holds.  If
\((u,p)\) is suitable in \(Q_1\), and if there is a strong predictor \(g\) with
\[
\|g\|_{L^q_tL^p_x(Q_1)}\le L,
\qquad
\frac2q+\frac3p=1,
\qquad
p>3,
\]
such that
\[
\|u-g\|_{X(Q_\theta)}
+\operatorname{PressTail}(Q_1,Q_\theta)
\le\varepsilon(L),
\]
then \(u\) is regular in a smaller cylinder \(Q_{\theta/2}\).  Here \(X\) is the
pressure-cleaned local energy-critical control norm used in the predictive
distance, and the pressure-tail term is measured in the same gauge used for the
relative Caffarelli-Kohn-Nirenberg gate.
\end{assumption}

\begin{assumption}[Predictor production and differentiability source theorem]
For every bounded incoming presentation there is a nonempty class of admissible
strong predictors.  The causal strong solution map \(\Phi\), from incoming data
to core predictors, is twice differentiable in the selected critical topology on
each bounded predictor class.  Quantitatively, for incoming data \(a\) and
perturbations \(h\),
\[
\|\Phi(a+h)-\Phi(a)-D\Phi(a)h\|_X
\le
C_L\|h\|_{\mathcal A}^2,
\]
with constants stable under parabolic rescaling and compatible with the selected
pressure-cleaning procedure.  The theorem also supplies realizable predictors
within the stated realization error of the relaxed minimax centre.
\end{assumption}

\begin{assumption}[Global residual representation source theorem]
Along every bad predictive chain, the registered local defects determine vectors
\(d_k\) in a fixed Hilbert defect space \(H_{\mathrm{def}}\).  There exist a
Hilbert space \(\mathcal K\), a residual vector \(\mathcal R\in\mathcal K\),
bounded maps \(A_k:\mathcal K\to H_{\mathrm{def}}\), and square-summable errors
\(\rho_k\) such that
\[
d_k=A_k\mathcal R+\rho_k.
\]
Moreover, the pulled-back innovation tests satisfy the Bessel bound required in
the moving-chart packing proposition.  The errors include, with no double
counting, pressure tails, cutoff losses, gauge changes, finite-window transport
errors, carrier recentering errors, and predictor-realization errors.
\end{assumption}

\begin{assumption}[Packet rigidity source theorem]
If a localized parabolic wave-packet family has quadratically small normal
leakage on a quantitatively connected interaction graph, then, after removing
pressure and gauge directions, it is close in the critical norm to a coherent
zero-cost branch.  The branch class contains the explicit Beltrami and
two-and-a-half-dimensional branches and every smooth Euler-coherent branch whose
projected quadratic nonlinearity is pure pressure in the selected gauge.  If the
packet family is not close to such a branch, the theorem produces a registered
normal defect detected by the Hilbert witness system.
\end{assumption}

\begin{assumption}[Carrier localization source theorem]
High-frequency carrier escape has one of three outcomes.  Either the carrier can
be recentered so that the dominant frequency returns to order one in a new
admissible presentation; or the escaped component is locally subcritical in every
carrier cell and therefore cannot maintain the bad-chain defect; or the escape
produces a registered normal, pressure, gauge, localization, or transport defect
which enters the global residual representation with square-summable error.
\end{assumption}

\begin{assumption}[Pressure and gauge synchronization source theorem]
For every scale block used in the iteration, pressure representatives, harmonic
tails, gauge normalizations, and mean-zero conventions can be chosen before the
block is evaluated.  Under those fixed choices, all pressure-force pairs used by
the Oseen estimate, the branch gate, and the relative
Caffarelli-Kohn-Nirenberg functional refer to the same pressure difference on
the core cylinder.  Any failure of synchronization is recorded as a pressure or
gauge defect in the global residual representation.
\end{assumption}

\begin{assumption}[Localization and critical collar source theorem]
For every controlled Oseen block, the fixed-profile derivative collars are
handled by the dyadic Hardy estimate proved above.  All remaining localization
terms--moving cutoffs, carrier collars, pressure tails created by localization,
and the critical collar sums
\[
\sum_k r_k^{-2}\int_{A_k}|w_k|^3,
\qquad
\sum_k r_k^{-2}\int_{A_k}|\pi_k^0|^{3/2}
\]
--either have a scale-invariant source bound, open a relative gate, produce a
registered pressure/localization/carrier defect, or enter a coherent branch
covered by the branch source theorem.
\end{assumption}

\begin{assumption}[Branch production and size source theorem]
Whenever the packet rigidity theorem selects a coherent zero-cost branch, the
presentation produces a strong branch predictor on a smaller cylinder together
with its synchronized pressure.  The predictor has a Serrin bound, its pressure
tail is controlled in the selected gauge, and its size data imply the relative
Caffarelli-Kohn-Nirenberg smallness needed by the branch gate, unless a
registered branch defect, pressure defect, force defect, or carrier escape is
produced.
\end{assumption}

\begin{assumption}[Finite presentation and update source theorem]
The iteration uses a finite presentation alphabet at each bounded complexity
level.  Predictor charts, pressure gauges, finite observational Gramians, branch
states, and carrier charts either stabilize along a tail of the bad chain, are
updated only finitely many times inside a fixed complexity level, or produce a
registered update/escape defect.  Thus an infinite bad chain cannot hide an
unpaid cost in infinitely many silent changes of presentation.
\end{assumption}

\begin{assumption}[Quantitative closure alternative source theorem]
For every \(L\) and \(\delta>0\), there are \(m=m(L,\delta)\) and
\(c=c(L,\delta)>0\) such that the following holds along any \(m\)-step controlled
block of a predictive chain.  If \(\mathfrak D_k\ge\delta\) throughout the block
and all registered normal, holonomy, pressure, carrier, update, and realization
defects in the block are at most \(c\), then either the defect decreases below
\(\delta/2\) at the next controlled scale, or the chain enters a coherent branch
or a carrier branch covered by the previous source theorems.
\end{assumption}

\subsection{Analytic work packages}

The source theorems above are the complete conditional interface used by the
proof.  Expanded into concrete analytic obligations, they require the following
statements.

\begin{enumerate}[label=(\roman*)]
\item \emph{Critical predictor construction.}  The incoming and lateral
observations must generate a nonempty family of strong predictors with Serrin
control, stable local pressure gauges, and quantitative dependence on the
incoming data.  The differentiability estimate must be proved in the same
critical topology used by the predictive defect, not in a stronger topology that
loses scale invariance.

\item \emph{Realization of minimax centres.}  The relaxed minimax centre exists
by Hilbert compactness, but the PDE iteration needs a realizable strong predictor
near that relaxed centre.  The realization theorem must either produce such a
predictor with the stated error or register the error as a summable realization
defect.

\item \emph{Global residual representation.}  Normal defects produced at
different scales must be represented as vectors \(d_k\) in a fixed defect
Hilbert space and must satisfy \(d_k=A_k\mathcal R+\rho_k\), with a Bessel bound
for the pulled-back innovation tests \(A_k^*e_k\).  The error vectors \(\rho_k\)
must control connection changes, cutoff changes, pressure gauges, carrier
recentering, finite-window truncation, and chart changes.  This is the statement
that makes the moving-chart packing proposition applicable to Navier-Stokes.

\item \emph{Packet pairing and row summability.}  The localized wave-packet
decomposition must supply actual pairings between physical residuals and
adjoint witnesses.  Repeated packet rows must obey either a Bessel/Carleson bound
or a direct operator bound in the critical topology.  A row that cannot be made
summable must exit by a gate, a coherent branch, carrier escape, or a registered
defect.

\item \emph{Packet rigidity.}  Small projected quadratic leakage must force the
localized packet family toward a coherent zero-cost branch after pressure and
gauge directions are removed.  If this rigidity conclusion fails, the failure
must produce a normal defect detected by the witness system.  The theorem must
cover localized packets, not only exact Fourier modes.

\item \emph{Pressure and gauge source closure.}  Pressure representatives,
harmonic tails, mean-zero conventions, gradient gauges, and pressure-force pairs
must be chosen in a block-independent way.  Variations of representatives,
pressure tails, gauge shrink losses, cutoff pressure losses, and pressure-force
couplings must be assigned once and only once to source terms, gates, or escape
alternatives.

\item \emph{Native force and old-data closure.}  Forcing terms created by
predictor updates, old visible data, finite-memory truncations, and transport of
past information must either be absorbed by the predictor construction or
registered as force defects with a summable source bound.  The proof cannot use
the same force contribution simultaneously as a predictor input and as a
separate residual source.

\item \emph{Cutoff and localization closure.}  The fixed-profile cutoff and
fixed-drift subchannels are proved above.  The remaining localization theorem
must handle moving cutoffs, annular critical mass, carrier collars, pressure
tails created by localization, and compatibility between localization cells and
the Oseen transfer.  Any failure must be routed to a localization defect,
carrier escape, pressure defect, or gate.

\item \emph{Oseen source closure beyond fixed subchannels.}  The local Oseen
energy transfer is proved, and fixed cutoff/drift costs are controlled.  What
remains is the source theorem for moving drifts, pressure-force pairs, native
forcing, localization-carrier-cutoff compatibility, and explicit Oseen error
terms.  These contributions must be separated from the packet, pressure, force,
and carrier source counts so that no cost is counted twice.

\item \emph{Carrier source closure.}  A high-frequency carrier must either be
recentred into a new admissible chart, become locally subcritical in each carrier
cell, or produce a summable registered defect.  The theorem must also control
carrier tails, containment of carrier cells in the physical cylinder, and
compatibility with the pressure gauge used by the final gate.

\item \emph{Coherent branch production.}  Once packet rigidity selects a
zero-cost branch, the proof needs an actual branch predictor with synchronized
pressure.  The branch theorem must produce outer branch data, interior size
data, derivative data when needed, pressure synchronization, and a route from
branch size to relative Caffarelli-Kohn-Nirenberg smallness.  If any of these
data fail, the failure must become a registered branch, pressure, force, or
carrier defect.

\item \emph{Branch radius and margin compatibility.}  The final gate is opened
on a smaller cylinder, so the branch radius, analytic radius, cutoff radius,
pressure tail allowance, and gate margin must be synchronized.  The theorem must
show that downstream consumption of these margins leaves a positive final gate
allowance, or else identifies a definite registered defect.

\item \emph{Nonlinear Caffarelli-Kohn-Nirenberg witnesses.}  Raw
Caffarelli-Kohn-Nirenberg mass is not a Hilbert coefficient.  Every use of such
mass must therefore be routed through the relative gate, converted to a genuine
linear source functional, controlled by an independent monotone/Carleson
estimate, or absorbed by compactness-rigidity.  Gated rows exit the iteration and
are not counted as source terms.

\item \emph{Finite presentation and update closure.}  The finite observational
atlas, predictor charts, pressure gauges, memory windows, coefficient ranks,
branch states, and carrier states must stabilize on controlled tails or pay a
finite update/escape cost.  This prevents an infinite bad chain from changing
presentations indefinitely without producing a summable source term or an exit.

\item \emph{Final floor and compactness closure.}  The no-gate condition, the
positive gate-excess defect, the compactness of bounded charts, and the
lower-semicontinuity of the final floor must be compatible with scale iteration.
This is the compactness statement that lets a persistent bad block survive in a
limit only if it retains a positive defect, and it is the point where the
quantitative closure alternative is consumed.
\end{enumerate}

\section{Conditional Newton Transfer}

Under the predictor production and differentiability source theorem, the Newton
bridge becomes a precise conditional proposition.

\begin{proposition}[Conditional saddle Newton bridge]
Assume critical predictor-map differentiability and the local Oseen transfer estimate. Let \(g_k\) be a relaxed minimax predictor at scale \(k\), and write \(u=g_k+w_k\). Then, after passing to the next core scale,
\[
\mathfrak D_{k+1}
\le
C_L\mathfrak D_k^2
+C_L
\left(
\mathfrak K_k^\perp
+\mathfrak H_k^\perp
+\mathfrak F_k^\perp
+\mathfrak V_k
\right)
+\eta_k.
\]
Here \(\mathfrak K_k^\perp\) is the non-absorbed normal quadratic defect, \(\mathfrak H_k^\perp\) the holonomy or gauge defect, \(\mathfrak F_k^\perp\) the pressure/localization defect, \(\mathfrak V_k\) the realization defect of the relaxed centre, and \(\eta_k\) the chosen minimization error.
\end{proposition}

\begin{proof}
The predictable linear part of \(w_k\) is absorbed by the first variation
\(D\Phi(a)h\). The tangential quadratic part is absorbed by the second-order
correction supplied by the differentiability source theorem. The non-absorbed
quadratic part is exactly \(\mathfrak K_k^\perp\); pressure, cutoff, and
connection errors give \(\mathfrak F_k^\perp\) and \(\mathfrak H_k^\perp\); and
non-realisability of the relaxed centre gives \(\mathfrak V_k\). Comparing the
suitable weak solution with the updated strong predictor, the local Oseen
estimate bounds the remainder on the smaller core. Since the updated predictor
is an admissible competitor for the minimax problem at the next scale, the
displayed inequality follows.
\end{proof}

\section{Modular Regularity Theorem}

\begin{theorem}[Conditional predictive regularity]
Let \((u,p)\) be a suitable weak solution near \(z_0\). Suppose that the
predictive presentation system at \(z_0\) satisfies the analytic source theorems
stated above, and that the explicit scale-selection errors \(\eta_k\) are
summable. Then there is no bad predictive chain at \(z_0\). Consequently, if the
predictive defect is the only obstruction to the predictive epsilon-regularity
hypothesis, then \(u\) is regular at \(z_0\).
\end{theorem}

\begin{proof}
Assume that a bad predictive chain exists. Then for some \(\delta_0>0\),
\[
\mathfrak D_k\ge\delta_0
\]
on all sufficiently small scales.

By the global residual representation source theorem and the Hilbertian packing
lemma, the innovative normal defects are square-summable. Hence, after discarding
finitely many initial scales, one can find arbitrarily long controlled blocks on
which the average registered normal defect is as small as required.  The pressure
and gauge synchronization theorem ensures that the pressure quantities appearing
in those blocks are the same quantities used by the Oseen estimate and by the
relative Caffarelli-Kohn-Nirenberg gate.

Apply the quantitative closure alternative on such a block. Since
\(\mathfrak D_k\ge\delta_0\) throughout the bad chain, the first outcome--defect
decay below \(\delta_0/2\)--is impossible on sufficiently late blocks. Therefore
the chain must enter either a coherent branch or a carrier branch, unless a
registered defect has been produced. A registered defect is also impossible
infinitely often with fixed positive size, because the global residual
representation and the packing lemma make the registered innovative cost
square-summable.

If the chain enters a Beltrami branch, the quadratic term is pressure-cleaned and
the Newton transfer gives decay of the predictive defect. If it enters a
two-and-a-half-dimensional branch, the reduced two-dimensional system and passive
scalar equation provide the strong predictor and again force decay, up to a
registered anisotropic defect. If it enters a more general Euler-coherent
zero-cost branch, the branch production and size source theorem supplies the
branch predictor and pressure data needed for the branch gate, unless a
registered branch, force, pressure, or carrier defect is produced. In the gate
case, the relative Caffarelli-Kohn-Nirenberg criterion opens the final gate and
the bad chain terminates. In the defect case, the defect has already been shown
to be square-summable and therefore cannot maintain the fixed lower bound
\(\delta_0\) indefinitely.

If the chain enters carrier escape, the carrier localization source theorem
either recenters the scale, makes the escaped contribution locally subcritical,
or produces one of the registered defects already controlled by packing.
Recentering gives a new admissible presentation and returns to the same
alternative; subcriticality contradicts persistence of the defect; and a
registered defect is square-summable. Every possible outcome contradicts the
existence of a bad chain with \(\mathfrak D_k\ge\delta_0\) on all sufficiently
small scales.

Therefore no bad predictive chain exists. The final implication is exactly the
predictive epsilon regularity source theorem.
\end{proof}

\section{Verified Components and Source Inputs}

The verified pieces are:
\begin{enumerate}[label=(\roman*)]
\item the Serrin-critical Oseen transfer estimate, with collar and pressure-tail terms explicitly paid for;
\item fixed-profile dyadic cutoff-collar control, fixed-drift collar stability,
and the pressure-pair source interface;
\item the Hilbert-space existence of relaxed minimax centres;
\item the abstract Bessel packing mechanism for innovative normal witnesses;
\item the exact two-mode null classification of the projected bilinear symbol;
\item exact absorption of Beltrami fields and exact reduction of two-and-a-half-dimensional fields;
\item the relative Caffarelli-Kohn-Nirenberg gate, conditional only on the stated predictor and pressure smallness;
\item compactness of the gate-excess defect under the selected chart topology;
\item suitable limit passage and normal-defect limit passage under the stated compactness hypotheses;
\item the separation between nonlinear Caffarelli-Kohn-Nirenberg witnesses and Hilbert innovations.
\end{enumerate}

The conditional theorem depends on the following scale-invariant PDE source
theorems:
\begin{enumerate}[label=(\roman*)]
\item production, differentiability, and realization of causal strong predictors;
\item construction of a global residual representation across a bad chain, with
a Bessel bound and square-summable pressure, cutoff, gauge, carrier, and
transport errors;
\item packet-level rigidity for localized nearly invisible quadratic
interactions;
\item carrier localization and recentering with pressure-tail control;
\item localization and critical velocity/pressure collar source bounds beyond
the fixed-profile Hardy subchannel;
\item pressure and gauge synchronization across Oseen, branch, and final-gate
quantities;
\item production of branch predictors, branch pressure data, and branch size
smallness;
\item finite presentation/update stabilization;
\item predictive epsilon regularity and the quantitative closure alternative.
\end{enumerate}

The useful conclusion of the paper is therefore a dependency graph with a proved
core. Presentation theory supplies the language of predictors, observables,
fibres, costs, and legal transfer; Oseen estimates supply the perturbative
transport; Hilbert packing supplies a square-summability mechanism for genuinely
new normal defects; relative Caffarelli-Kohn-Nirenberg gates convert predictor
smallness into regularity; and null geometry isolates the branches where the
quadratic term can be invisible.  The remaining analytic task is to prove the
listed source theorems in the concrete Navier-Stokes topology.

\begin{thebibliography}{99}

\bibitem[ABP]{ABP}
D. Albritton, T. Barker, and C. Prange,
\emph{Epsilon regularity for the Navier-Stokes equations via weak-strong uniqueness},
arXiv:2211.16188.

\bibitem[CKN]{CKN}
L. Caffarelli, R. Kohn, and L. Nirenberg,
\emph{Partial regularity of suitable weak solutions of the Navier-Stokes equations},
Communications on Pure and Applied Mathematics 35 (1982), 771-831.

\bibitem[ESS]{ESS}
L. Escauriaza, G. Seregin, and V. Sverak,
\emph{\(L_{3,\infty}\)-solutions of Navier-Stokes equations and backward uniqueness},
Russian Mathematical Surveys 58 (2003), 211-250.

\bibitem[Leray]{Leray}
J. Leray,
\emph{Sur le mouvement d'un liquide visqueux emplissant l'espace},
Acta Mathematica 63 (1934), 193-248.

\bibitem[Serrin]{Serrin}
J. Serrin,
\emph{On the interior regularity of weak solutions of the Navier-Stokes equations},
Archive for Rational Mechanics and Analysis 9 (1962), 187-195.

\bibitem[Tao]{Tao}
T. Tao,
\emph{Quantitative bounds for critically bounded solutions to the Navier-Stokes equations},
preprint, 2019.

\bibitem[Wolf]{Wolf}
J. Wolf,
\emph{On the local pressure of the Navier-Stokes equations and related systems},
Advances in Differential Equations 17 (2012), 305-338.

\end{thebibliography}

\end{document}
